\scanpage{9}
On suppose maintenant pour simplifier que
l'Anneau $\uline{O}_S$ est \unsure{\uline{strict\textsuperscript{mt}}} \uline{local}
\inserted{(p.ex. anneau \unsure{faisceau dans} top. ét)}. On dispose alors
\margin{\unsure{On supp.} de plus $2$ inv. dans $\uline{O}_S$}
\[
WG(S)\xrightarrow{w}H^*(S,\mathbb F_2)
\]
(classes de Stiefel Whitney), caractérisé (pour
$S$ variable) par

a) \uline{fonctorialité} en $S$

b) $\unsure{w(M,q)}$ dans le cas où $M$ de rang $1$
(normalisation)

c) $w(x+y)=w(x)w(y)$
\[
w(\unsure{x|m})-1\in H^+(S,\mathbb F_2)\qquad\text{i.e. }w^0(x)=1.
\]

\margin{\deleted{\unsure{par splitting} \unsure{on se ramène à} \unsure{prouver}
$w(\alpha(E))=1$, $E$ de rang $1$.}}
\deleted{Un calcul facile montre que $w\circ\alpha=1$, donc
$w$ se factorise en}
\deleted{\[
W(S)\xrightarrow{w}1+\widehat H^{*+}(S,\mathbb F_2)
\]}

\deleted{\uline{NB} $w(M,q)$ \illegible{3} \unsure{que}}

\uline{NB} Les arguments habituels de splitting
montrent que $w(M\otimes N)$ et $w(\lambda^{\unsure{i}}M)$ se
calculent par les formules \inserted{universelles} habituelles en
termes de $w(M),w(N)$ et $\rho(M),\rho(N)$ (rangs
des Modules $M,N$). Donc on a un \inserted{$\lambda$-}hom
\[
WG(S)\longrightarrow H^0(S,\mathbb Z)\times\deleted{\unsure{H}}(1+\widehat H^{*+}(S,\mathbb F_2))
\]

\mspage{8}
